% ====================================================================
% Chapter 1: Graphite anode ICA tail theory
% Scope: discharge/delithiation branch; thermodynamic stage + kinetic tail.
% ====================================================================
\documentclass[11pt,a4paper]{article}
\usepackage[margin=22mm]{geometry}
\usepackage{kotex}
\usepackage{amsmath,amssymb,mathtools,bm}
\usepackage{booktabs,longtable,array}
\usepackage{enumitem}
\usepackage{amsthm}
\usepackage{hyperref}
\usepackage{fancyhdr}
\usepackage{lastpage}
\usepackage{setspace}
\setstretch{1.13}
\setlength{\parskip}{0.45em}\setlength{\parindent}{0pt}\setlength{\headheight}{14pt}
\hypersetup{colorlinks=true, linkcolor=blue, urlcolor=blue, citecolor=blue,
  pdftitle={Graphite anode ICA tail theory - Chapter 1}}
\pagestyle{fancy}\fancyhf{}
\lhead{Graphite anode ICA tail theory --- Ch.1}\rhead{\thepage/\pageref{LastPage}}
\renewcommand{\headrulewidth}{0.4pt}
\newtheorem*{groundingbox}{Grounding}
\newtheorem*{boundbox}{Validity bound}
\newtheorem*{flagbox}{Flagged assumption}
\newtheorem*{keybox}{Keystone}
\newcommand{\dd}{\mathrm{d}}
\newcommand{\eff}{\mathrm{eff}}
\newcommand{\eq}{\mathrm{eq}}
\newcommand{\app}{\mathrm{app}}
\newcommand{\drive}{\mathrm{drive}}
\newcommand{\bg}{\mathrm{bg}}
\newcommand{\cell}{\mathrm{cell}}
\newcommand{\ext}{\mathrm{ext}}
\newcommand{\OCV}{\mathrm{OCV}}
\newcommand{\tot}{\mathrm{tot}}
\newcommand{\tail}{\mathrm{tail}}
\newcommand{\AL}[1]{\textsf{[AL-#1]}}
\title{\textbf{리튬이온전지 흑연 음극 ICA($dQ/dV$) tail 이론 --- Chapter 1}\\[0.4em]
\large Thermodynamic stage, electrode-potential-assisted relaxation barrier, and first-pass fitting approximation}
\author{}\date{}
\begin{document}
\maketitle

\section*{서: 관찰에서 fitting approximation까지}
\addcontentsline{toc}{section}{서: 관찰에서 fitting approximation까지}
본 장은 다음 관찰에서 출발한다.
\begin{quote}\itshape
흑연 음극의 ICA(incremental capacity, $dQ/dV$) peak에서 상전이에 의해 생기는 tail은 온도가 낮을수록 길게 늘어지고,
온도가 높을수록 상대적으로 짧게 끝난다. 또한 현재 전위 상태와 C-rate가 바뀌면 peak overlap과 tail 모양도 달라진다.
\end{quote}

이 현상은 equilibrium thermodynamics만으로는 충분히 설명되지 않는다. Ideal lattice-gas equilibrium peak의 thermal width는
$RT/F$ scale이므로, 다른 효과가 없다면 낮은 온도에서 더 좁아지는 방향을 준다. 따라서 관측되는 low-temperature long tail의
주된 원인은 equilibrium peak width 자체가 아니라, transition progress가 equilibrium target을 유한한 속도로 따라가는
kinetic lag에서 찾아야 한다. 본 장의 목적은 이 kinetic lag를 electrode potential이 보조하는 effective relaxation barrier와
연결하고, 마지막에는 실제 ICA 데이터에서 바로 쓸 수 있는 simple fitting approximation을 도출하는 것이다.

논리 흐름은 다음 여섯 단계다.
\begin{enumerate}[label=\textbf{(\arabic*)},leftmargin=2.2em]
\item Charge conservation이 internal anode potential $V_n$을 정한다. $V_n$은 lookup value가 아니라 internal state와 external charge가 함께 결정하는 implicit variable이다.
\item Equilibrium target $\xi_{\eq,j}(V_n,T)$는 충분히 느린 scan에서 transition이 도달할 목표를 준다.
\item 실제 progress $\xi_j$는 finite rate $k_j$로 target을 따라간다. $k_j$가 작으면 tail이 길어진다.
\item Electrode-potential drive $\mathcal A_j$는 effective relaxation barrier를 낮추는 bounded approximation으로 들어가며, 그 결과 $k_j$와 tail length $L$이 바뀐다.
\item Barrier distribution은 exponential map을 통해 relaxation-length spectrum이 되고, 관측 tail은 single-mode kernel의 중첩으로 표현된다.
\item Dominant single-mode limit에서는 $Y_{\tail}(q)\approx B\exp[-(q-q_a)/L]$가 되어 Chapter 1만으로 first-pass fitting이 가능하다.
\end{enumerate}

\begin{groundingbox}
모든 equation은 먼저 어떤 physical quantity를 계산하는지 말로 정의한 뒤 제시하고, 뒤에서 각 항의 meaning과 unit을 해석한다.
Assumption은 \S\ref{sec:convention}의 ledger에서 established, bounded, flagged로 분리한다. Paper 또는 textbook 근거가
확인되지 않은 항목은 theory identity처럼 쓰지 않는다. 본 장은 discontinuous switch, 임의 cap, 임의 threshold를 physical model로 쓰지 않는다.
\end{groundingbox}

\tableofcontents
\newpage

\section{기호와 단위}\label{sec:notation}
진행률 기호는 $\xi$로 통일한다. 일부 문헌의 $\theta$는 occupation/progress variable로만 언급하고 본문 기호로 쓰지 않는다.
본 장은 discharge 중 graphite anode가 delithiation되는 branch를 다룬다. External discharge coordinate는
$q=Q_{\ext}/Q_{\cell}$이며 증가 방향으로 둔다.

\begin{longtable}{p{0.18\textwidth} p{0.12\textwidth} p{0.62\textwidth}}
\toprule 기호 & 단위 & 의미 \\ \midrule \endhead
$q$ & --- & normalized external discharge coordinate, $q=Q_{\ext}/Q_{\cell}$ \\
$Q_{\ext},Q_{\cell}$ & C & external passed charge와 reference cell capacity \\
$I$ & A & applied current. 크기는 $|I|$, discharge sign은 별도 $s_I$ \\
$V_n$ & V & internal graphite anode potential. charge conservation의 solution \\
$V_{n,\app}$ & V & measured apparent anode potential, $V_n+s_I|I|R_n$ \\
$V_{n,\drive}$ & V & transition rate에 들어가는 drive potential. Chapter 1 기본형은 $V_{n,\drive}=V_n$ \\
$V_{n,\OCV}$ & V & equilibrium limit에서 charge conservation이 주는 anode OCV \\
$R_n$ & $\Omega$ & apparent-axis shift를 주는 effective polarization coefficient \\
$j$ & --- & discharge curve에서 분리 가능한 effective transition label \\
$\xi_j,\xi_{\eq,j}$ & --- & branch-local transition progress와 equilibrium target, $0\le\xi\le1$ \\
$U_j,w_j$ & V & effective transition center potential과 equilibrium width \\
$Q_{j,\tot}$ & C & transition $j$가 discharge capacity에 주는 positive contribution \\
$Q_{\bg},C_{\bg}$ & C, C/V & unresolved background capacity와 $C_{\bg}=\partial Q_{\bg}/\partial V_n$ \\
$\mathcal A_j$ & J/mol & electrode-potential drive, $s_{\phi,j}F(V_{n,\drive}-U_j)$ \\
$s_{\phi,j}$ & --- & discharge branch에서 $V_{n,\drive}>U_j$가 progress를 촉진하도록 정한 sign factor \\
$G$ or $g$ & J/mol & potential assistance 전 intrinsic relaxation barrier value \\
$\Delta G_{\eff,j}$ & J/mol & effective relaxation barrier, $G-\chi_j\mathcal A_j$ \\
$\chi_j$ & --- & Level-A relaxation-barrier sensitivity. Chapter 1에서는 Butler--Volmer transfer coefficient $\beta_j$와 동일시하지 않음 \\
$k_j,k_0$ & 1/s & relaxation rate와 prefactor \\
$L$ & --- & $q$ coordinate의 relaxation length, $L=|I|/(Q_{\cell}k_j)$ \\
$L_\varphi$ & V & voltage-axis tail length, local approximation $L_\varphi\simeq |\dd V/\dd q|_{q_a}L$ \\
$\rho_G(G;T)$ & mol/J & intrinsic barrier probability density \\
$p_L(L)$ & 1 & relaxation-length probability density. $p_L(L)\dd L$은 mode fraction \\
$A_L(L)$ & 1 & relaxation-length amplitude spectrum. $A_L(L)\dd L$은 $[L,L+\dd L]$ mode들이 싣는 tail progress amplitude \\
$a(L)$ & --- & probability spectrum을 amplitude spectrum으로 바꾸는 residual/capacity/accessibility factor \\
$\Theta,\Theta_{\eq}$ & --- & capacity-weighted aggregate progress와 equilibrium target \\
$\Theta_{\tail},\Theta_0$ & --- & post-peak tail progress와 $q_a$에서 남은 single-mode residual amplitude \\
$Q_p$ & C & aggregate progress의 capacity scale. 보통 local transition group에 대해 $Q_p=\sum_jQ_{j,\tot}$ \\
$q_a,x$ & ---, --- & tail-window start와 local coordinate $x=q-q_a$ \\
$Y_{\tail}$ & measured ICA unit & first-pass fitting에 쓰는 baseline-subtracted ICA tail signal \\
\bottomrule
\end{longtable}

전류 $I$는 $\mathrm{C/s}$이고 $Q_{\cell}$은 C이다. Constant-current segment에서는
\begin{equation}
Q_{\ext}(t)=\int_0^t|I|\dd t',\qquad q=\frac{Q_{\ext}}{Q_{\cell}},\qquad
\frac{\dd q}{\dd t}=\frac{|I|}{Q_{\cell}}.
\label{eq:q_time}
\end{equation}

\section{Convention과 assumption ledger}\label{sec:convention}
\textbf{Branch convention.} Graphite lithiation literature에서는 stage evolution을 lithiation order로 말하는 경우가 많다.
그러나 본 장은 measured discharge/delithiation branch를 기준으로 transition을 label한다. 따라서 $j$는 crystallographic
stage number가 아니라 branch-local peak label이다. 각 $j$에 대해 $\xi_j=0$은 해당 discharge peak가 아직 진행되지 않은 상태,
$\xi_j=1$은 그 branch-local transition이 지난 상태를 뜻한다.

\textbf{Barrier convention.} $G$는 intrinsic relaxation barrier다. $\Delta G_{\eff}=G-\chi\mathcal A$는 Butler--Volmer 의미의
microscopic forward barrier가 아니다. Chapter 1에서는 $\xi_j$가 $\xi_{\eq,j}$로 접근하는 scalar relaxation mobility를 정하는
Level-A effective barrier로만 쓴다. Directional forward/backward barrier splitting은 Chapter 3의 reaction-kinetic extension에서
다룬다.

\begin{longtable}{p{0.10\textwidth} p{0.16\textwidth} p{0.30\textwidth} p{0.34\textwidth}}
\toprule ID & status & source type & allowed use \\ \midrule \endhead
\AL{1} & established & charge conservation in porous-electrode theory \cite{doyle1993,newman} & $V_n$은 charge balance로 결정된다. \\
\AL{2} & established/bounded & lattice-gas and regular-solution thermodynamics \cite{mckinnon1983,bazant2013} & ideal logistic은 유도식, nonideal effective width는 bounded coarse-graining. \\
\AL{3} & established & graphite staging observations \cite{dahn1991,ohzuku1993,funabiki1999ea,funabiki1999jes} & staging transition이 ICA peak를 만든다. \\
\AL{4} & bounded & low-rate OCV/GITT calibration & $U_j,w_j,Q_{j,\tot},Q_{\bg}$는 equilibrium-stage fitting quantity. \\
\AL{5} & validity gate & positive differential capacity / implicit-function condition \cite{newman,bazant2013} & fitting window에서 $C_{\bg}>0$ 필요. \\
\AL{6} & established & ohmic and transport polarization separation \cite{newman,bardfaulkner} & $V_{n,\app}$와 $V_n$을 분리해야 $\chi_j$를 해석할 수 있다. \\
\AL{7} & mathematical gate & range and monotonicity condition & charge-balance solution은 admissible window에서만 존재. \\
\AL{10} & established & transition-state theory \cite{eyring1935} & rate scale은 $\exp[-\Delta G^\ddagger/(RT)]$를 따른다. \\
\AL{11} & bounded & linear free-energy/BEP/BV analogy \cite{evanspolanyi1938,bardfaulkner,bazant2013} & transition 근방의 scalar relaxation barrier에 대한 first-order potential assistance. \\
\AL{12} & established/bounded & Onsager relaxation and two-state kinetic algebra \cite{onsager1931,degrootmazur1962} & local first-order relaxation toward a target. \\
\AL{13} & flagged/bounded & heterogeneous electrode microstructure plus graphite cooperative transition literature & barrier-only independent mode distribution은 mean-field approximation. \\
\AL{14} & established mechanism, bounded application & relaxation-time distributions \cite{svare2000,johnston2006,lindsey1980} & barrier distribution이 broad relaxation spectrum을 만들 수 있음. graphite ICA 적용은 model contribution. \\
\AL{15} & established bound & Marcus curvature \cite{marcus1956} & linear barrier assistance는 small-drive approximation. \\
\AL{16} & bounded analogy & Fredholm ratio-substitution methods \cite{lee2011,son2013,jcp2017} & unknown-ratio closure idea만 사용. Fredholm formula를 Volterra problem에 그대로 이식하지 않음. \\
\bottomrule
\end{longtable}

\section{흑연 staging과 effective transition}\label{sec:staging}
Graphite lithiation/delithiation은 staged intercalation structure를 지나며, stage transformation은 뚜렷한 ICA peak를 남긴다
\cite{dahn1991,ohzuku1993,funabiki1999ea,funabiki1999jes}. 실제 measured $dQ/dV$에서는 여러 microscopic stage transformation이
겹칠 수 있으므로, 본 장은 모든 measured peak를 crystallographic stage boundary와 일대일 대응시키지 않고 effective transition
$j=1,\dots,N_p$로 다룬다.

하나의 effective transition 내부에도 local mode가 존재할 수 있다. Particle size, defect density, grain boundary, stress,
local accessibility는 relaxation barrier를 바꿀 수 있다. Chapter 1에서는 우선 이 heterogeneity를 intrinsic barrier $G$의
분포로 취급하고, effective transition의 equilibrium target은 공유된다고 둔다. 이것이 \emph{barrier-only disorder} approximation이다.
이는 실제 graphite phase-boundary motion이 microscopic하게 독립이라는 뜻이 아니라, tail 해석을 위해 barrier heterogeneity의 효과만
분리해 보는 reduced model이다.

\section{무대 I: charge conservation과 internal potential}\label{sec:charge}
External current는 지나간 charge를 정하지만, internal anode potential을 직접 정하지 않는다. Internal potential은 stored charge가
passed charge와 같아지도록 정해지는 값이다:
\begin{equation}
\boxed{\;Q_{\cell}q=Q_{\bg}(V_n,T)+\sum_{j=1}^{N_p}Q_{j,\tot}\xi_j.\;}
\label{eq:charge_balance}
\end{equation}
좌변은 experiment에서 아는 값이다. 우변은 internal potential과 transition progress에 의존한다. 따라서 주어진
$(q,T,\{\xi_j\})$에서 Eq.~\eqref{eq:charge_balance}를 만족하는 값이 $V_n$이다.

$Q_{\bg}$는 그림에서 임의로 빼는 visual baseline이 아니다. 선택한 effective transition으로 분리하지 않은 capacity contribution이다.
그 differential capacity는
\begin{equation}
C_{\bg}(V_n,T)=\frac{\partial Q_{\bg}}{\partial V_n}
\label{eq:Cbg}
\end{equation}
이고, local fitting window에서는
\begin{equation}
C_{\bg}(V_n,T)>0
\label{eq:monotone}
\end{equation}
이어야 한다. 이는 numerical patch가 아니라 implicit solution이 local하게 well-posed이기 위한 validity condition이다.

Equilibrium에서는 $\xi_j=\xi_{\eq,j}(V_n,T)$이고, 같은 charge-balance equation이 equilibrium anode potential
$V_{n,\OCV}(q,T)$를 준다. 즉 OCV는 charge conservation의 equilibrium special solution이지, charge conservation을 대체하는
별도 lookup이 아니다.

Measured apparent potential은 polarization을 포함한다:
\begin{equation}
V_{n,\app}=V_n+s_I|I|R_n(q,T,|I|).
\label{eq:Vapp}
\end{equation}
따라서 measured data를 internal axis로 옮기려면 charge-balance consistency와 polarization correction이 모두 필요하다.
First-pass fit에서는 $R_n$을 low-rate data, pulse, current-interrupt, impedance 같은 독립 정보로 먼저 제한한 뒤 tail 변화를
$\chi_j$로 해석해야 한다.

\section{무대 II: equilibrium target}\label{sec:eq}
Minimal equilibrium target은 lattice-gas 또는 regular-solution picture에서 얻을 수 있다. Occupation/progress를 $\xi$라 하면
chemical potential은
\begin{equation}
\mu(\xi)=\mu^0+RT\ln\frac{\xi}{1-\xi}+\Omega_j(1-2\xi)
\label{eq:mu_latticegas}
\end{equation}
로 쓸 수 있다. Logarithmic term은 configurational entropy에서 온다. $N$개 site 중 $n=N\xi$개가 occupied이면
$W=N!/[n!(N-n)!]$이고, Stirling approximation을 쓰면 site 1 mol 기준으로
\begin{equation}
S_{\mathrm{mix}}=-R[\xi\ln\xi+(1-\xi)\ln(1-\xi)]
\end{equation}
이다. $G_{\mathrm{mix}}=-TS_{\mathrm{mix}}$를 $\xi$로 미분하면 $RT\ln[\xi/(1-\xi)]$가 나온다. Regular-solution enthalpy
$\Omega_j\xi(1-\xi)$를 미분하면 $\Omega_j(1-2\xi)$가 된다.

Ideal limit $\Omega_j=0$에서 electrode potential과 equilibrium을 맞추면
\begin{equation}
RT\ln\frac{\xi_{\eq,j}}{1-\xi_{\eq,j}}=s_{\phi,j}F(V_n-U_j)
\label{eq:eq_balance}
\end{equation}
이다. Standard-state chemical potential은 $U_j$에 흡수했다. $z=s_{\phi,j}(V_n-U_j)/w_j$이고 ideal case에서
$w_j=RT/F$라 두면
\begin{equation}
\boxed{\;\xi_{\eq,j}(V_n,T)=\frac{1}{1+\exp[-s_{\phi,j}(V_n-U_j)/w_j]}\;}
\label{eq:logistic}
\end{equation}
가 된다. $\Omega_j\ne0$, particle heterogeneity, finite plateau width가 중요하면 $w_j$는 $RT/F$에 강제로 묶지 않고
low-rate equilibrium data로 calibration하는 effective width가 된다.

Ideal target의 derivative는
\begin{equation}
\frac{\partial\xi_{\eq,j}}{\partial V_n}=\frac{s_{\phi,j}}{w_j}\xi_{\eq,j}(1-\xi_{\eq,j})
\label{eq:dxidV}
\end{equation}
이다. 따라서 ideal thermal contribution만 보면 낮은 $T$에서 $RT/F$가 작아져 peak가 좁아진다. 이것이 실제 graphite
equilibrium peak가 항상 정확히 $RT/F$로 좁아진다는 뜻은 아니다. Graphite staging은 cooperative transition 성격이 있으므로
low-rate OCV/GITT로 equilibrium shape를 확인해야 한다. 다만 관측된 low-temperature long tail을 pure ideal thermal broadening만으로
설명하는 것은 논리적으로 충분하지 않다.

\section{주역 I: electrode-potential-assisted relaxation barrier}\label{sec:rate}
Kinetic variable은 equilibrium target 자체가 아니라 actual progress가 그 target으로 접근하는 rate다. Local electrode-potential drive를
\begin{equation}
\mathcal A_j=s_{\phi,j}F(V_{n,\drive}-U_j)
\label{eq:Aj}
\end{equation}
로 정의한다. 단위는 J/mol이다. 본 장의 discharge branch convention에서는 $\mathcal A_j>0$일 때 selected discharge transition의
progress가 촉진되도록 $s_{\phi,j}$를 잡는다.

Transition-state theory에서는 relaxation rate가 Boltzmann factor에 비례한다:
\begin{equation}
k_j\propto \exp\left[-\frac{\Delta G_j^\ddagger}{RT}\right].
\end{equation}
Transition center 근방에서는 scalar relaxation barrier에 대해 bounded linear free-energy approximation을 쓴다:
\begin{equation}
\boxed{\;\Delta G_{\eff,j}=G-\chi_j\mathcal A_j.\;}
\label{eq:Geff}
\end{equation}
여기서 $G$는 local mode의 intrinsic relaxation barrier이고, $\chi_j$는 electrode-potential drive가 \emph{relaxation mobility}를
얼마나 바꾸는지 나타내는 coefficient다. 이 $\chi_j$는 아직 Butler--Volmer transfer coefficient $\beta_j$가 아니다.
Butler--Volmer/Marcus theory는 보통 forward barrier와 reverse barrier를 서로 다르게 바꾸어 forward/reverse ratio를 바꾼다.
반면 Chapter 1의 Level A에서는 equilibrium target은 thermodynamics가 정하고, $\Delta G_{\eff}$는 그 target으로 접근하는 속도만
바꾼다. Directional kinetics는 Chapter 3에서 추가한다.

Eyring-type prefactor를 쓰면
\begin{equation}
\boxed{\;k_j(T,V_{n,\drive})=k_0(T)
\exp\left[-\frac{G-\chi_j\mathcal A_j}{RT}\right],\qquad
k_0(T)=\frac{k_BT}{h}\kappa(T)\exp\left(\frac{\Delta S_a}{R}\right).\;}
\label{eq:kact}
\end{equation}
이다. Prefactor의 묶음 방식은 달라질 수 있지만, tail shape에 중요한 것은 $k_j$가 $G-\chi_j\mathcal A_j$에 exponential하게
민감하다는 점이다. $|\mathcal A_j|$가 reorganization/curvature scale과 비슷해지는 strong-drive 영역에서는 linear barrier
approximation을 global law로 쓰지 않는다.

\begin{keybox}
Two-state kinetic algebra는 first-order relaxation form의 의미를 분명히 해준다. Local forward/reverse rate를 $r_j^+$와
$r_j^-$라 하면
\begin{equation}
\dot\xi_j=r_j^+(1-\xi_j)-r_j^-\xi_j=(r_j^++r_j^-)
\left[\frac{r_j^+}{r_j^++r_j^-}-\xi_j\right].
\end{equation}
따라서 local window에서 rate를 상수로 보면 $k_j=r_j^++r_j^-$이고
$\xi_{\eq,j}=r_j^+/(r_j^++r_j^-)$다. Chapter 1은 이 $k_j$를 scalar relaxation rate로 쓴다. 즉
$\chi_j\mathcal A_j$를 Chapter 3의 microscopic forward barrier lowering이라고 주장하지 않는다.
\end{keybox}

다른 physical bottleneck이 relaxation을 제한하면 lumped first-order approximation으로
\begin{equation}
\frac{1}{k_j^{\eff}}=\frac{1}{k_j}+\frac{1}{k_{j,\lim}}
\label{eq:kphys}
\end{equation}
을 둘 수 있다. 이는 transport, finite-site, diffusion, accessibility bottleneck을 하나의 reduced time scale로 묶는 bounded model이다.
독립 근거 없이 free fitting knob으로 추가하면 안 된다.

\section{주역 II: progress dynamics와 single-mode kernel}\label{sec:dyn}
Constant-current segment에서 local mode의 first-order relaxation은
\begin{equation}
\frac{\dd\xi_j}{\dd t}=k_j[\xi_{\eq,j}(V_n,T)-\xi_j]
\label{eq:dxidt}
\end{equation}
이다. Eq.~\eqref{eq:q_time}을 이용하면
\begin{equation}
\frac{\dd\xi_j}{\dd q}=\frac{Q_{\cell}k_j}{|I|}[\xi_{\eq,j}-\xi_j]
\label{eq:dxidq}
\end{equation}
가 된다.

Post-peak window의 시작점을 $q_a$라 하고 residual을 $r=\xi_{\eq,j}-\xi_j$로 둔다. 이 window에서 equilibrium target이 거의
flat하여 $\dd\xi_{\eq,j}/\dd q\simeq0$이고, $\mathcal A_j$ 변화가 작아 $k_j$가 거의 constant이면
\begin{equation}
\frac{\dd r}{\dd q}=-\frac{Q_{\cell}k_j}{|I|}r=-\frac{1}{L}r,
\qquad
\boxed{\;L=\frac{|I|}{Q_{\cell}k_j}.\;}
\label{eq:Ldef}
\end{equation}
이다. 따라서
\begin{equation}
r(q)=r(q_a)\exp\left[-\frac{q-q_a}{L}\right]
\label{eq:residual_decay}
\end{equation}
이고, observed progress contribution은 residual 자체가 아니라 derivative이므로
\begin{equation}
\boxed{\;\frac{\dd\xi_j}{\dd q}=\frac{r(q_a)}{L}\exp\left[-\frac{q-q_a}{L}\right].\;}
\label{eq:single_dxidq}
\end{equation}
이다. 이 $1/L$ factor가 뒤의 kernel factor의 출처다. $L$은 normalized charge coordinate $q$에서의 length이므로 dimensionless다.

온도 경향도 여기서 나온다. $T$가 낮아지고 $G-\chi\mathcal A$가 positive로 유지되면 $k_j$가 작아지고 $L$이 커진다.
큰 $L$은 $q$축에서 더 천천히 decay하는 tail, 즉 longer tail을 의미한다.

\section{주역 III: barrier distribution에서 relaxation-length spectrum으로}\label{sec:spectrum}
실제 electrode는 하나의 intrinsic barrier만 갖지 않는다. Intrinsic relaxation barrier의 probability density를
$\rho_G(G;T)$라 하자. Barrier-only disorder approximation에서는 서로 다른 $G$를 가진 mode들이 같은 branch-local
equilibrium target을 공유하지만 서로 다른 rate로 relax한다고 둔다:
\begin{equation}
\xi_j(t)=\int_0^\infty \rho_G(g;T)\xi_j(g,t)\dd g,
\qquad
\frac{\dd\xi_j(g,t)}{\dd t}=k_j(g)[\xi_{\eq,j}-\xi_j(g,t)].
\label{eq:xi_dist}
\end{equation}
이 approximation은 local defect가 $U_j,w_j$, capacity까지 바꿀 가능성을 무시한다. 그런 효과가 dominant이면 model을 확장해야 한다.
Chapter 1에서는 barrier heterogeneity가 tail에 주는 효과만 분리하기 위해 이 reduced form을 사용한다.

Eq.~\eqref{eq:kact}를 Eq.~\eqref{eq:Ldef}에 넣으면
\begin{equation}
\boxed{\;L(G,T,V_{n,\drive})=\frac{|I|}{Q_{\cell}k_0(T)}
\exp\left[\frac{G-\chi_j\mathcal A_j}{RT}\right].\;}
\label{eq:L_of_G}
\end{equation}
이다. 이것이 temperature amplification의 핵심이다. 고정된 $G$ spread는 $\ln L$에서 $1/(RT)$에 비례해 확대된다.
따라서 low temperature에서는 relaxation-length spectrum이 넓어지고, long-lived tail mode의 영향이 커진다.

Length coordinate의 probability density는 mode fraction 보존으로 얻는다:
\begin{equation}
p_L(L)\dd L=\rho_G(G)\dd G.
\end{equation}
역변환은 $G(L)=\chi_j\mathcal A_j+RT\ln(L/L_0)$, $L_0=|I|/(Q_{\cell}k_0)$이고,
\begin{equation}
\frac{\dd G}{\dd L}=\frac{RT}{L}
\end{equation}
이므로
\begin{equation}
\boxed{\;p_L(L)=\rho_G(G(L);T)\frac{RT}{L},\qquad
L\ge L_{\min}=L_0\exp\left[-\frac{\chi_j\mathcal A_j}{RT}\right].\;}
\label{eq:spectrum_prob}
\end{equation}
이다. Lower limit은 artificial switch가 아니라 $G\ge0$에 해당하는 domain condition이다.

Observed tail에는 probability spectrum만이 아니라 amplitude spectrum이 필요하다. $a(L)$에 $q_a$에서 남은 residual amplitude,
local capacity/accessibility, 독립적으로 제한된 weight를 묶으면
\begin{equation}
\boxed{\;A_L(L)=a(L)p_L(L)=a(L)\rho_G(G(L);T)\frac{RT}{L},\qquad L\ge L_{\min}.\;}
\label{eq:spectrum}
\end{equation}
이다. $A_L(L)\dd L$은 relaxation length가 $[L,L+\dd L]$인 mode들이 post-peak progress amplitude에 싣는 양이다.
일반적으로 normalized density가 아니다.

\section{관측 tail: kernel integral}\label{sec:kernel}
각 length mode는 Eq.~\eqref{eq:single_dxidq}의 derivative form을 갖는다. 그 derivative contribution을 length axis에서 합하면
\begin{equation}
\boxed{\;\frac{\dd\Theta_{\tail}}{\dd q}(q)
\simeq \int_{L_{\min}}^\infty A_L(L)\frac{1}{L}
\exp\left[-\frac{q-q_a}{L}\right]\dd L,
\qquad q\ge q_a.\;}
\label{eq:kernel_integral}
\end{equation}
이다. Integration variable은 relaxation length이고, result는 external charge coordinate에서의 progress derivative다.
$A_L$ 안의 $RT/L$은 $G\to L$ coordinate transformation에서 왔고, kernel의 $1/L$은 residual decay를 progress derivative로
바꿀 때 왔다. 두 factor의 origin은 다르다.

Single mode가 dominant이면
\begin{equation}
A_L(L)=\Theta_0\delta(L-L_*)
\end{equation}
이고, Eq.~\eqref{eq:kernel_integral}은
\begin{equation}
\frac{\dd\Theta_{\tail}}{\dd q}\simeq\frac{\Theta_0}{L_*}
\exp\left[-\frac{q-q_a}{L_*}\right]
\label{eq:single_from_kernel}
\end{equation}
로 줄어든다. 단순히 $A_L=\delta(L-L_*)$라고만 쓰면 amplitude가 빠진다. 따라서 amplitude-bearing delta form이 필요하다.

\begin{boundbox}
Barrier-distribution relaxation과 stretched relaxation은 broader relaxation literature에서 확립된 mechanism이다
\cite{svare2000,johnston2006,lindsey1980}. 그러나 이를 graphite ICA tail에 적용하는 것은 본 model의 bounded claim이다.
따라서 본 장은 $\rho_G$를 one tail curve에서 uniquely invert하지 않고, 독립 정보 또는 constrained forward model로 다룬다.
\end{boundbox}

\section{Self-consistency: feedback이 들어가는 위치}\label{sec:volterra}
$\Theta$는 stored charge에 기여하므로 Eq.~\eqref{eq:charge_balance}를 통해 $V_n$을 바꾼다. $V_n$이 바뀌면 equilibrium target도
바뀐다. 이 feedback은 실제로 존재하지만, first-pass exponential fitting formula를 얻기 위한 필수 조건은 아니다. 여기서는 이후
correction이 어디에 들어가는지만 명확히 둔다.

Single mode에 대해
\begin{equation}
\frac{\dd\Theta}{\dd q}=\frac{1}{L}[\Theta_{\eq}(q)-\Theta(q)]
\end{equation}
라고 쓰면, integrating factor $e^{q/L}$를 곱해 적분하여
\begin{equation}
\Theta(q)=\Theta(0)e^{-q/L}+\int_0^q\frac{1}{L}e^{-(q-q')/L}\Theta_{\eq}(q')\dd q'
\label{eq:singlemode_integral}
\end{equation}
를 얻는다. Spectrum과 charge-balance feedback을 넣으면 causal Volterra equation
\begin{equation}
\Theta(q)=\Theta_{\mathrm{init}}(q)+\int_0^q\mathcal K(q,q')
\Theta_{\eq}\left(V_n(q',\Theta(q')),T\right)\dd q'
\label{eq:volterra}
\end{equation}
가 된다. 여기서
\begin{equation}
\mathcal K(q,q')=\int_{L_{\min}}^\infty A_L(L)\frac{1}{L}e^{-(q-q')/L}\dd L.
\end{equation}
Kernel은 target에 곱해진다. 이 integral stage에서 kernel을 다시 residual $[\Theta_{\eq}-\Theta]$에 곱하면 relaxation을 두 번 세는
오류가 된다.

\subsection{Plan B: reference evaluation}\label{sec:planB}
Plan B는 Chapter 1의 main theory가 아니라 reference calculation이다. Barrier grid $g_m$를 두고 각 local ODE를 적분하면서,
각 time 또는 charge step마다 Eq.~\eqref{eq:charge_balance}로 $V_n$을 푼다. Spectrum이 넓거나 feedback이 강할 때 Plan B는
Plan A 또는 simple approximation을 검증하는 기준해 역할을 한다.

\subsection{Plan A: optional analytic closure}\label{sec:planA}
사용자의 JCP paper와 refs.~6,7은 Fredholm integral equation of the second kind를 풀기 위해 unknown-ratio approximation을 쓴다
\cite{lee2011,son2013,jcp2017}. 2017 paper에서 Eq.~(32)는 Fredholm equation이고, Eq.~(34)는 unknown solution ratio를 reference
ratio로 대체한다. 현재 graphite equation은 Volterra equation이므로 그 formula를 그대로 가져오지 않는다. Controlled ratio
substitution이라는 idea만 가져온다.

Target을 baseline solution $\Theta^{(0)}$ 주변에서 linearization하면
\begin{equation}
\Theta_{\eq}(V_n(\Theta))\approx a(q)+b(q)\Theta(q),
\qquad
b(q)=\frac{\partial\Theta_{\eq}}{\partial V_n}\frac{\partial V_n}{\partial\Theta}
=-\frac{Q_p}{C_{\bg}}\frac{\partial\Theta_{\eq}}{\partial V_n}
\label{eq:selfcoupling}
\end{equation}
이다. $b\le0$는 $\partial\Theta_{\eq}/\partial V_n\ge0$이고 $C_{\bg}>0$인 window에서만 성립한다. Ratio approximation
$\Theta(q')/\Theta(q)\approx\Theta^{(0)}(q')/\Theta^{(0)}(q)$를 쓰면
\begin{equation}
\boxed{\;\Theta_A(q)=\frac{c(q)}{1-\dfrac{1}{\Theta^{(0)}(q)}
\displaystyle\int_0^q\mathcal K(q,q')b(q')\Theta^{(0)}(q')\dd q'}\;,}
\label{eq:closed}
\end{equation}
이고 $c(q)=\Theta_{\mathrm{init}}(q)+\int_0^q\mathcal K(q,q')a(q')\dd q'$이다. 이 closure는 bounded approximation이므로
quantitative fitting expression으로 쓰기 전에 Plan B와 비교해야 한다.

\section{ICA/DVA relation과 simple fitting formula}\label{sec:fiteq}
Charge conservation을 $q$로 미분하면
\begin{equation}
Q_{\cell}=C_{\bg}\frac{\dd V_n}{\dd q}+\sum_jQ_{j,\tot}\frac{\dd\xi_j}{\dd q}
\end{equation}
이다. $Q_p\Theta=\sum_jQ_{j,\tot}\xi_j$로 정의하면
\begin{equation}
\dd Q=C_{\bg}\dd V_n+Q_p\dd\Theta
\label{eq:dQsplit}
\end{equation}
이다. 또한 $\dd Q=Q_{\cell}\dd q$이므로
\begin{equation}
\boxed{\;\frac{\dd Q}{\dd V_n}=\frac{C_{\bg}(V_n,T)}{1-Q_p(\dd\Theta/\dd Q)}
=\frac{C_{\bg}(V_n,T)}{1-(Q_p/Q_{\cell})(\dd\Theta/\dd q)}.\;}
\label{eq:fiteq}
\end{equation}
따라서 $\dd\Theta/\dd q$의 tail은 ICA tail로 나타난다. 그러나 $\dd\Theta/\dd q$가 exponential이라고 해서 measured ICA가
정확히 exponential이라는 뜻은 아니다. Exact reduced expression은 tail contribution에 대해 rational form이고, exponential fitting은
small-tail 및 locally constant baseline approximation에서만 성립한다.

Apparent axis에서는
\begin{equation}
\frac{\dd V_{n,\app}}{\dd q}=\frac{\dd V_n}{\dd q}+s_I|I|\frac{\partial R_n}{\partial q},
\qquad
\frac{\dd Q_{\ext}}{\dd V_{n,\app}}=\frac{Q_{\cell}}{\dd V_{n,\app}/\dd q}
\label{eq:app_axis}
\end{equation}
이다. 따라서 voltage-axis tail length는 $L_\varphi\simeq |\dd V_{n,\app}/\dd q|_{q_a}L$이다. 이를
$|\dd V_n/\dd q|_{q_a}L$로 바꾸는 것은 chosen tail window에서 $\partial R_n/\partial q$가 무시 가능할 때만 허용된다.

\subsection{Simple real-data fitting approximation}\label{sec:ch1_simplefit}
$x=q-q_a$라 하자. Single dominant mode tail에서는
\begin{equation}
\frac{\dd\Theta_{\tail}}{\dd q}\simeq\frac{\Theta_0}{L}e^{-x/L},
\qquad
\boxed{\;L=\frac{|I|}{Q_{\cell}k_0(T)}
\exp\left[\frac{G-\chi_j\mathcal A_j}{RT}\right].\;}
\label{eq:simplefit}
\end{equation}
이다. Baseline-subtracted ICA tail signal을
\begin{equation}
Y_{\tail}(q)=\left(\frac{\dd Q}{\dd V}\right)_{\mathrm{meas}}
-\left(\frac{\dd Q}{\dd V}\right)_{\mathrm{baseline}}
\label{eq:Ytail_def}
\end{equation}
로 정의한다. Baseline은 low-rate equilibrium fit, peak 밖에서 제한된 local polynomial/spline, 또는 Chapter 1의 charge-balance
background에서 얻는 slowly varying contribution이다. Tail을 흡수하도록 자유롭게 흔드는 curve가 아니다.

$(Q_p/Q_{\cell})(\dd\Theta/\dd q)\ll1$이고, chosen tail window에서 voltage-axis conversion이 거의 constant이면
Eq.~\eqref{eq:fiteq}는
\begin{equation}
\boxed{\;Y_{\tail}(q)\approx B\exp[-(q-q_a)/L].\;}
\label{eq:Ytail_simple}
\end{equation}
로 줄어든다. $B$에는 $C_{\bg}$, $Q_p/Q_{\cell}$, $\Theta_0/L$, 선택한 voltage-axis conversion이 들어간다. Tail shape의
핵심 parameter는 $L$이다.

\textbf{First-pass fitting protocol.}
\begin{enumerate}[label=\textbf{F\arabic*.},leftmargin=2.4em]
\item Low-rate OCV/GITT로 equilibrium stage, 즉 $U_j,w_j,Q_{j,\tot},Q_{\bg}$를 먼저 고정한다.
\item $R_n$은 $\chi_j$보다 먼저 독립적으로 제한한다. 둘을 같은 tail에서 동시에 자유 fitting하면 식별성이 무너진다.
\item Peak maximum 이후, equilibrium derivative가 measurement noise floor보다 작거나 semi-log residual이 locally linear한 구간을 tail window로 잡는다. 이 규칙은 data noise에 묶어야 하며 universal arbitrary threshold가 아니다.
\item $q_a$는 fitting 전에 고정한다. $q_a$를 free parameter로 두면 $L$과 강하게 correlated된다.
\item 선택한 window에서 $\ln Y_{\tail}=\ln B-(q-q_a)/L$를 회귀하여 slope로 $L$을 얻는다.
\item $L$은
\begin{equation}
G-\chi_j\mathcal A_j=RT\ln\left(\frac{Q_{\cell}k_0(T)L}{|I|}\right)
\label{eq:barrier_from_L}
\end{equation}
로 effective barrier와 연결된다. $k_0(T)$가 독립적으로 알려지지 않으면 single curve의 absolute value보다 multi-temperature 또는 multi-potential difference를 사용한다.
\item Arrhenius analysis에서는 Eyring prefactor와 entropy term을 고려한다. $\ln(1/L)$ 대 $1/T$의 단순 slope가
$-(\Delta H_a-\chi\mathcal A)/R$를 뜻하려면 $\mathcal A$가 고정되고 prefactor variation이 보정되거나 negligible해야 한다.
\item Fixed temperature에서 local potential-dependence는
\begin{equation}
\frac{\partial\ln L}{\partial V_{n,\drive}}=-\frac{\chi_js_{\phi,j}F}{RT}
\label{eq:chi_slope}
\end{equation}
로 $\chi_j$를 제한한다. 이 식은 small-drive window에서만 쓴다.
\end{enumerate}

Semi-log tail에 systematic curvature가 있으면 single-mode approximation이 실패한 것이다. 이는 단순 fitting failure가 아니라
relaxation-length spectrum, changing $L(q)$, 또는 self-consistency feedback의 신호다.

\section{Temperature, potential, and C-rate effects}\label{sec:falsify}
Reduced formula는 세 contribution을 분리한다.
\begin{longtable}{p{0.25\textwidth} p{0.30\textwidth} p{0.35\textwidth}}
\toprule layer & terms & role \\ \midrule \endhead
Equilibrium stage & $U_j(T),w_j(T),Q_{\bg}(V_n,T)$ & peak location과 equilibrium target \\
Kinetic lag & $k_j(T,V_{n,\drive}),\rho_G(G;T)$ & tail length와 spectrum \\
Apparent-axis polarization & $R_n(q,T,|I|)$ & measured voltage-axis shift와 distortion \\
\bottomrule
\end{longtable}

Temperature는 $L$의 exponential factor를 바꾼다. 낮은 $T$에서는 $G/(RT)$가 커지고 $\ln L$ distribution이 넓어지므로,
kinetic-lag model이 맞다면 low temperature에서 long tail이 자연스럽게 나타난다. Electrode potential은
\begin{equation}
\frac{\partial\ln L}{\partial V_{n,\drive}}=-\frac{\chi_js_{\phi,j}F}{RT}
\end{equation}
를 통해 $L$을 바꾼다. 본 discharge sign convention에서는 potential assistance가 $k_j$를 키워 tail을 짧게 만드는 방향을 준다.

C-rate effect는 단일 방향으로 단정하면 안 된다. $k_j$가 고정이면 $L=|I|/(Q_{\cell}k_j)$이므로 $q$ coordinate에서 current 증가는
$L$을 키운다. 반대로 current 증가가 $V_{n,\drive}$를 shift하여 $k_j$를 크게 만들면 $L$은 작아진다. Apparent polarization은
measured voltage axis에서 peak 위치와 overlap도 바꾼다. 따라서 C-rate trend는 current prefactor, potential-assisted rate increase,
voltage-axis polarization의 competition으로 해석해야 한다.

Model은 다음 기준으로 falsifiable하다.
\begin{enumerate}[label=\textbf{N\arabic*.},leftmargin=2.4em]
\item Low-rate equilibrium data가 chosen equilibrium target shape와 $w_j$를 지지해야 한다.
\item Extracted $L$의 Arrhenius-type behavior는 $\mathcal A_j$와 prefactor variation을 통제한 뒤 평가해야 한다.
\item Local slope $\partial\ln L/\partial V_{n,\drive}$는 small-drive window에서 Eq.~\eqref{eq:chi_slope}의 sign과 scale을 따라야 한다.
\item $\rho_G$는 independent physical information 또는 low-dimensional forward model로 제한해야 한다. Single tail curve만으로 uniquely determine되지 않는다.
\item $R_n$과 $\chi_j$를 한 dataset에서 동시에 free parameter로 두면 covariance 때문에 physical interpretation이 무너질 수 있다.
\end{enumerate}

\section{부록: self-consistency evaluation principle}\label{sec:ch1_numeric}
Single-mode approximation이 실패하면 Plan B를 reference calculation으로 쓴다. Experimental tail window에서 보이는 relaxation length를
덮도록 barrier grid $\{g_m\}$를 잡는다. 각 $g_m$에 대해
\begin{equation}
\frac{\dd\xi_j(g_m)}{\dd t}=k_j(g_m)[\xi_{\eq,j}(V_n,T)-\xi_j(g_m)]
\end{equation}
를 적분하고, 각 step에서 charge conservation으로 $V_n$을 푼다. Algebraic constraint는 local charge-balance derivative condition이
만족될 때 well posed이다. Plan A는 같은 temperature, rate, tail window에서 Plan B와의 차이가 measurement noise 및 discretization floor보다
작을 때만 받아들인다.

\section*{Chapter 2--5로의 전달}
Chapter 1은 $V_n$, $V_{n,\app}$, $V_{n,\drive}$, $\xi_j$, $\xi_{\eq,j}$, $k_j$, $L$, $A_L$의 convention을 고정한다.
Chapter 2는 이 equilibrium stage 위에서 reversible heat와 entropy coefficient를 다루고, Chapter 3은 Level-A scalar relaxation을
directional forward/backward kinetics로 확장한다. Chapter 4는 heat terms를 double counting 없이 결합하고, Chapter 5는 Chapter 1
convention을 유지한 채 charge/discharge branch와 hysteresis를 추가해야 한다.

\begin{thebibliography}{99}
\bibitem{doyle1993} M. Doyle, T. F. Fuller, J. Newman, ``Modeling of Galvanostatic Charge and Discharge of the Lithium/Polymer/Insertion Cell,'' \emph{J. Electrochem. Soc.} \textbf{140}, 1526 (1993). DOI: 10.1149/1.2221597.
\bibitem{newman} J. Newman, K. E. Thomas-Alyea, \emph{Electrochemical Systems}, 3rd ed., Wiley (2004). ISBN 978-0-471-47756-3.
\bibitem{mckinnon1983} W. R. McKinnon, R. R. Haering, ``Physical Mechanisms of Intercalation,'' in \emph{Modern Aspects of Electrochemistry} No.~15, Plenum, New York (1983), p.~235.
\bibitem{bazant2013} M. Z. Bazant, ``Theory of Chemical Kinetics and Charge Transfer Based on Nonequilibrium Thermodynamics,'' \emph{Acc. Chem. Res.} \textbf{46}, 1144 (2013). DOI: 10.1021/ar300145c.
\bibitem{bardfaulkner} A. J. Bard, L. R. Faulkner, \emph{Electrochemical Methods}, 2nd ed., Wiley (2001). ISBN 978-0-471-04372-0.
\bibitem{eyring1935} H. Eyring, ``The Activated Complex in Chemical Reactions,'' \emph{J. Chem. Phys.} \textbf{3}, 107 (1935). DOI: 10.1063/1.1749604.
\bibitem{evanspolanyi1938} M. G. Evans, M. Polanyi, ``Inertia and Driving Force of Chemical Reactions,'' \emph{Trans. Faraday Soc.} \textbf{34}, 11 (1938). DOI: 10.1039/TF9383400011.
\bibitem{marcus1956} R. A. Marcus, ``On the Theory of Oxidation-Reduction Reactions Involving Electron Transfer. I,'' \emph{J. Chem. Phys.} \textbf{24}, 966 (1956). DOI: 10.1063/1.1742723.
\bibitem{onsager1931} L. Onsager, ``Reciprocal Relations in Irreversible Processes. I,'' \emph{Phys. Rev.} \textbf{37}, 405 (1931). DOI: 10.1103/PhysRev.37.405.
\bibitem{degrootmazur1962} S. R. de Groot, P. Mazur, \emph{Non-Equilibrium Thermodynamics}, North-Holland (1962); Dover (1984).
\bibitem{dahn1991} J. R. Dahn, ``Phase Diagram of Li$_x$C$_6$,'' \emph{Phys. Rev. B} \textbf{44}, 9170 (1991). DOI: 10.1103/PhysRevB.44.9170.
\bibitem{ohzuku1993} T. Ohzuku, Y. Iwakoshi, K. Sawai, ``Formation of Lithium-Graphite Intercalation Compounds in Nonaqueous Electrolytes,'' \emph{J. Electrochem. Soc.} \textbf{140}, 2490 (1993). DOI: 10.1149/1.2220849.
\bibitem{funabiki1999ea} A. Funabiki, M. Inaba, T. Abe, Z. Ogumi, ``Nucleation and Phase-Boundary Movement upon Stage Transformation in Lithium-Graphite Intercalation Compounds,'' \emph{Electrochim. Acta} \textbf{45}, 865 (1999). DOI: 10.1016/S0013-4686(99)00290-X.
\bibitem{funabiki1999jes} A. Funabiki, M. Inaba, T. Abe, Z. Ogumi, ``Stage Transformation of Lithium-Graphite Intercalation Compounds Caused by Electrochemical Lithium Intercalation,'' \emph{J. Electrochem. Soc.} \textbf{146}, 2443 (1999). DOI: 10.1149/1.1391953.
\bibitem{baessler1993} H. B\"assler, ``Charge Transport in Disordered Organic Photoconductors,'' \emph{Phys. Status Solidi B} \textbf{175}, 15 (1993). DOI: 10.1002/pssb.2221750102.
\bibitem{svare2000} I. Svare, S. W. Martin, F. Borsa, ``Stretched Exponentials with $T$-Dependent Exponents from Fixed Distributions of Energy Barriers for Relaxation Times in Fast-Ion Conductors,'' \emph{Phys. Rev. B} \textbf{61}, 228 (2000). DOI: 10.1103/PhysRevB.61.228.
\bibitem{johnston2006} D. C. Johnston, ``Stretched Exponential Relaxation Arising from a Continuous Sum of Exponential Decays,'' \emph{Phys. Rev. B} \textbf{74}, 184430 (2006). DOI: 10.1103/PhysRevB.74.184430.
\bibitem{lindsey1980} C. P. Lindsey, G. D. Patterson, ``Detailed Comparison of the Williams-Watts and Cole-Davidson Functions,'' \emph{J. Chem. Phys.} \textbf{73}, 3348 (1980). DOI: 10.1063/1.440530.
\bibitem{jcp2017} K. Lee, S. Lee, C. H. Choi, S. Lee, ``Effects of External Electric Field and Anisotropic Long-Range Reactivity on Charge Separation Probability,'' \emph{J. Chem. Phys.} \textbf{147}, 144111 (2017). DOI: 10.1063/1.5000882.
\bibitem{lee2011} S. Lee, C. Y. Son, J. Sung, S.-H. Chong, ``Communication: Propagator for Diffusive Dynamics of an Interacting Molecular Pair,'' \emph{J. Chem. Phys.} \textbf{134}, 121102 (2011). DOI: 10.1063/1.3565476.
\bibitem{son2013} C. Y. Son, J. Kim, J.-H. Kim, J. S. Kim, S. Lee, ``An Accurate Expression for the Rates of Diffusion-Influenced Bimolecular Reactions with Long-Range Reactivity,'' \emph{J. Chem. Phys.} \textbf{138}, 164123 (2013). DOI: 10.1063/1.4802584.
\end{thebibliography}

\end{document}
